
















\subsection{可对角化性}
前一小节表明，
虽然相似的矩阵必定是矩阵等价的，但
反之不成立。
有些矩阵等价类会分裂成两个或更多
相似类；例如，非奇异的 $\nbyn{2}$ 矩阵
构成一个矩阵等价类，但不止一个相似类。

下面的图示说明了这一点。
实线曲线表示矩阵等价类，
而虚线分界线标出相似类。
每颗星都是表示其所在相似类的一个矩阵。
%  so each
% similarity class has one.
我们不能把用于矩阵等价的规范形式，
即分块部分单位矩阵，
用作相似（关系）的规范形式，
因为每个矩阵等价类中只有一个
部分单位矩阵。
\begin{center}
  \includegraphics{jc/mp/ch5.5}
\end{center}
要为各相似类的
代表元设计一种规范形式，
我们自然会在此前工作的基础上展开。
于是，如果某个相似类中确实含有部分单位矩阵，
那么它就应当表示该类。
除此之外，代表元应当尽可能简单。
%(the partial identities are simple in that they consist mostly of zeros).

部分单位形式最简单的推广是对角形式。

\begin{definition} \label{df:Diagonalizable}
%<*df:Diagonalizable>
若一个变换\definend{可对角化}\index{diagonalizable|(}%
\index{transformation!diagonalizable}，
则它在定义域与陪域取同一组基时
具有对角的表示。
一个\definend{可对角化矩阵}\index{matrix!diagonalizable}
是指与某个对角矩阵相似的矩阵：~\( T \) 可对角化
当且仅当存在非奇异的 \( P \) 使得 \( PTP^{-1} \) 是对角的。
%</df:Diagonalizable>
\end{definition}

\begin{example} \label{ex:DiagTwoByTwo}
矩阵
\begin{equation*}
  \begin{mat}[r]
     4 &-2 \\
     1 &1
  \end{mat}
\end{equation*}
是可对角化的。
\begin{equation*}
  \begin{mat}[r]
     2  &0   \\
     0  &3
  \end{mat}
  =
  \begin{mat}[r]
    -1  &2  \\
     1  &-1
  \end{mat}
  \begin{mat}[r]
     4  &-2 \\
     1  &1
  \end{mat}
  \begin{mat}[r]
    -1  &2  \\
     1  &-1
  \end{mat}^{-1}
\end{equation*}
\end{example}

下面我们将看到如何求出矩阵~$P$，但首先我们指出，
并非每个矩阵都与对角矩阵相似，
所以对角形式不足以充当
相似（关系）的规范形式。

\begin{example}  \label{ex:MatrixNotDiagonalizable}
%<*ex:NotDiagonalizable>
这个矩阵不是可对角化的。
\begin{equation*}
  N=\begin{mat}[r]
       0  &0  \\
       1  &0
    \end{mat}
\end{equation*}
$N$ 不是零矩阵这一事实意味着
它不可能与
零矩阵相似，因为零矩阵只与它自身相似。
于是，如果 $N$ 与某个对角矩阵~$D$ 相似，
那么 $D$ 在对角线上至少有一个非零元。

关键之处在于 $N$ 的某个幂是零矩阵，具体地，
$N^2$ 是零矩阵。
这意味着，对于任意由 \( N \)
关于某组 $B,B$ 所表示的映射~$n$，
复合 \( \composed{n}{n} \) 是零映射。
这又意味着，
表示~$n$ 的任意
关于某组 $\hat{B},\hat{B}$ 的矩阵，其平方都是零矩阵。
但对任意非零对角矩阵~$D^2$，$D^2$ 的元素
是 $D$ 的元素的平方，
所以 $D^2$ 不可能是零矩阵。
因此 $N$ 不是可对角化的。
%</ex:NotDiagonalizable>
\end{example}

所以并非每个相似类都含有一个对角矩阵。
下面我们来刻画矩阵何时可对角化。
% However, some similarity classes do contain a diagonal matrix and
% the canonical form that we are developing has the property that if
% a matrix can be diagonalized then the diagonal matrix is the canonical
% representative of its similarity class.

\begin{lemma} \label{lm:DiagIffBasisOfEigens}
%<*lm:DiagIffBasisOfEigens>
变换 \( t \) 可对角化当且仅当
存在一组基
\( B=\sequence{\vec{\beta}_1,\ldots,\vec{\beta}_n } \)
和标量 \( \lambda_1,\ldots,\lambda_n \)，使得
\( t(\vec{\beta}_i)=\lambda_i\vec{\beta}_i \)
对每个 \( i \) 都成立。
%</lm:DiagIffBasisOfEigens>
\end{lemma}

\begin{proof}
%<*pf:DiagIffBasisOfEigens>
考虑一个对角的表示矩阵。
\begin{equation*}
   \rep{t}{B,B}=
   \begin{pmat}{c@{\hspace*{1em}}c@{\hspace*{1em}}c}
      \vdots                    &       &\vdots                     \\
      \rep{t(\vec{\beta}_1)}{B} &\cdots &\rep{t(\vec{\beta}_n)}{B}  \\
      \vdots                    &       &\vdots
   \end{pmat}
   =
   \begin{pmat}{c@{\hspace*{1em}}c@{\hspace*{1em}}c}
      \lambda_1   &       &0         \\
      \vdots      &\ddots &\vdots    \\
      0           &       &\lambda_n
   \end{pmat}
\end{equation*}
考虑该基的一个成员
关于基 $\rep{\vec{\beta}_i}{B}$ 的表示。
对角矩阵与表示向量的乘积
\begin{equation*}
   \rep{t(\vec{\beta}_i)}{B}
   =\begin{pmat}{c@{\hspace*{1em}}c@{\hspace*{1em}}c}
      \lambda_1   &       &0         \\
      \vdots      &\ddots &\vdots    \\
      0           &       &\lambda_n
   \end{pmat}
  \colvec{0 \\ \vdots \\ 1 \\ \vdots \\ 0}
  =\colvec{0 \\ \vdots \\ \lambda_i \\ \vdots \\ 0}
\end{equation*}
便具有
所述的作用。
%</pf:DiagIffBasisOfEigens>
\end{proof}

\begin{example}     \label{ex:DiagUpperTrian}
要对角化
\begin{equation*}
   T=\begin{mat}[r]
      3  &2  \\
      0  &1
   \end{mat}
\end{equation*}
我们把 $T$ 取为一个变换关于
标准基的表示 $\rep{t}{\stdbasis_2,\stdbasis_2}$，并寻找一组基
\( B=\sequence{\vec{\beta}_1,\vec{\beta}_2} \)，使得
\begin{equation*}
  \rep{t}{B,B}
  =
  \begin{mat}
    \lambda_1  &0          \\
    0          &\lambda_2
  \end{mat}
\end{equation*}
也就是说，使得
$t(\vec{\beta}_1)=\lambda_1\vec{\beta}_1$
以及 $t(\vec{\beta}_2)=\lambda_2\vec{\beta}_2$。
\begin{equation*}
  \begin{mat}[r]
     3  &2  \\
     0  &1
  \end{mat}
  \vec{\beta}_1=\lambda_1\cdot\vec{\beta}_1
  \qquad
  \begin{mat}[r]
     3  &2  \\
     0  &1
  \end{mat}
  \vec{\beta}_2=\lambda_2\cdot\vec{\beta}_2
\end{equation*}
我们要找的是使方程
\begin{equation*}
 \begin{mat}[r]
    3  &2  \\
    0  &1
 \end{mat}
 \colvec{b_1 \\ b_2}=x\cdot\colvec{b_1 \\ b_2}
\end{equation*}
有不全为 $0$ 的解 $b_1$ 与 $b_2$ 的标量 \( x \)
（零向量不是任何基的成员）。
这就是一个线性方程组。
\begin{equation*}
  \begin{linsys}{2}
     (3-x)\cdot b_1  &+  &2\cdot b_2       &=  &0  \\
                     &   &(1-x)\cdot b_2   &=  &0
  \end{linsys}
\tag*{($*$)}\end{equation*}
先关注最下面的方程。
有两种情形：要么
$b_2=0$，要么~$x=1$。

在 \( b_2=0 \) 的情形，
第一个方程给出：要么 $b_1=0$，要么 \( x=3 \)。
既然我们已经排除了 $b_2=0$ 且~$b_1=0$ 同时成立的可能性，
就只剩下第一个对角元 $\lambda_1=3$。
于是，($*$) 的第一个方程是 $0\cdot b_1+2\cdot b_2=0$，
因此与 $\lambda_1=3$ 相关联的是
第二分量为零而第一分量自由的向量。
\begin{equation*}
     \begin{mat}[r]
        3  &2  \\
        0  &1
     \end{mat}
     \colvec{b_1 \\ 0}=3\cdot\colvec{b_1 \\ 0}
\end{equation*}
为得到第一个基向量，
任取一个非零的 $b_1$ 即可。
\begin{equation*}
   \vec{\beta}_1=\colvec[r]{1 \\ 0}
\end{equation*}

($*$) 最下面方程的另一种情形是 $\lambda_2=1$。
这时 ($*$) 的第一个方程是 $2\cdot b_1+2\cdot b_2=0$，因此
与这种情形相关联的是
第二分量为第一分量之相反数的向量。
\begin{equation*}
     \begin{mat}[r]
        3  &2  \\
        0  &1
     \end{mat}
     \colvec{b_1 \\ -b_1}=1\cdot\colvec{b_1 \\ -b_1}
\end{equation*}
由此得到第二个基向量：
从这些向量中任取一个非零的即可。
\begin{equation*}
   \vec{\beta}_2=\colvec[r]{1 \\ -1}
\end{equation*}

现在画出相似图
（回忆一下，我们使用的标量是复数，
而非实数），
\begin{equation*}
     \begin{CD}
            \C^2_{\wrt{\stdbasis_2}}
               @>t>T>
               \C^2_{\wrt{\stdbasis_2}}       \\
            @V{\scriptstyle\identity} VV
               @V{\scriptstyle\identity} VV \\
            \C^2_{\wrt{B}}
               @>t>D>
               \C^2_{\wrt{B}}
      \end{CD}
\end{equation*}
并注意到矩阵 $\rep{\identity}{B,\stdbasis_2}$ 很容易求出，
由此得到这个对角化。
\begin{equation*}
   \begin{mat}[r]
     3  &0  \\
     0  &1
   \end{mat}
   =
   \begin{mat}[r]
     1  &1  \\
     0  &-1
   \end{mat}^{-1}
   \begin{mat}[r]
     3  &2  \\
     0  &1
   \end{mat}
   \begin{mat}[r]
     1  &1  \\
     0  &-1
   \end{mat}
\end{equation*}
\end{example}

本节余下部分将围绕那个例子展开，进一步考察
\nearbylemma{lm:DiagIffBasisOfEigens} 的性质，
并给出一种
求这些 $\lambda$ 的简捷方法。
下一节则围绕 \nearbyexample{ex:MatrixNotDiagonalizable} 展开，
以理解是什么可能阻碍对角化。
再下一节将这两者结合起来，给出一种
在某种意义上尽可能简单的规范形式。

